
%--------------------------------------------------------------------------
% 核心包加载（适用于 pdflatex）
%--------------------------------------------------------------------------
% 基础功能包
\usepackage{graphicx}                             % 图形插入支持
\usepackage{longtable}                            % 长表格支持
\usepackage{amsmath,amsthm,amssymb,amsfonts}      % 数学符号和字体支持
\usepackage{thmtools}                             % 高级定理工具包（提供declaretheorem）
\usepackage[dvipsnames]{xcolor}                   % 颜色支持
\usepackage{tikz}                                 % 矢量绘图
\usepackage{tikz-cd}                              % 交换图绘制
\usepackage{mathtools}                            % 数学工具包（提供多种数学环境和命令）
\usepackage{mathrsfs}                             % 花体字母支持
\usepackage{bbm}                                  % 黑板粗体字母支持

% 页面设置和布局
\usepackage{geometry}
\geometry{
	a4paper,
	top=2.5cm,
	bottom=2.5cm,
	left=2.5cm,
	right=2.5cm,
}
\renewcommand{\baselinestretch}{1.3}

%--------------------------------------------------------------------------
% Article-style section headings and table of contents
%--------------------------------------------------------------------------
% Follow the standard LaTeX article hierarchy: prominent section headings,
% bold section entries in the contents, and indented subsection entries.
% \makeatletter
% \renewcommand\section{\@startsection{section}{1}{\z@}%
% 	{-3.5ex \@plus -1ex \@minus -.2ex}%
% 	{2.3ex \@plus .2ex}%
% 	{\normalfont\Large\bfseries\def\@secnumfont{\bfseries}}}
% \renewcommand*\l@section{\@tocline{1}{1em \@plus\p@}{0pt}{1.5em}{\bfseries}}
% \renewcommand*\l@subsection{\@tocline{2}{0pt}{1.5em}{2.3em}{}}
% \makeatother


\usepackage[backend=biber,style=alphabetic,backref]{biblatex} % 参考文献管理
\ExecuteBibliographyOptions{maxbibnames=99}                       % 显示所有作者
\ExecuteBibliographyOptions{maxalphanames=4,minalphanames=4}     % 标签显示全部作者（如 DLOZ22）
\usepackage{xurl} % allow long URLs to break anywhere

% 交互功能
\usepackage[
	colorlinks,
	linkcolor=NavyBlue,
	urlcolor=magenta,
	citecolor=cyan,
	hypertexnames=false,
]{hyperref}       % 超链接支持
\usepackage[nameinlink]{cleveref} % 智能引用（自动添加前缀）


% 列表环境标号配置
\usepackage{enumitem}       % Required for custom labels
\setlist[enumerate,1]{label=(\alph*)} % 第一层级：小写字母加括号 (a), (b), ...)
\setlist[enumerate,2]{label=(\roman*)} % 第二层级：小写罗马数字加括号 (i), (ii), ...)

%--------------------------------------------------------------------------
% 定理环境可视化配置
%--------------------------------------------------------------------------
% 定义具体定理环境（共享计数器，正文使用斜体）
\theoremstyle{plain}
\newtheorem{definition}{Definition}[section]
\declaretheorem[sibling=definition,name=Proposition]{proposition}
\declaretheorem[sibling=definition,name=Theorem]{theorem}
\declaretheorem[sibling=definition,name=Lemma]{lemma}
\declaretheorem[sibling=definition,name=Corollary]{corollary}
\declaretheorem[sibling=definition,name=Conjecture]{conjecture}
\declaretheorem[sibling=definition,name=Question]{question}
\declaretheorem[sibling=definition,name=Claim]{claim}
\declaretheorem[numbered=yes, name=Theorem]{mainthm} % main theorem with global numbering
\renewcommand{\themainthm}{\Alph{mainthm}} % Display as A, B, C, ...
\theoremstyle{remark}
\declaretheorem[sibling=definition,name=Remark]{remark}
\declaretheorem[sibling=definition,name=Example]{example}


% 定义step子环境（在proof环境中使用，每个proof重置计数器）
\newcounter{step}
\newcounter{globalstep} % Global unique step counter for anchors
\newenvironment{step}[1][]{%
	\refstepcounter{step}%
	\refstepcounter{globalstep}% This creates unique anchors
	\par\smallskip%
	\noindent\textbf{Step \thestep\if\relax\detokenize{#1}\relax\else\ (#1)\fi.}\ %
}{%
	\par\smallskip%
}
\renewcommand{\theglobalstep}{\arabic{step}}

% 定义case子环境（在proof环境中使用，每个proof重置计数器）
\newcounter{case}
\newcounter{globalcase} % Global unique case counter for anchors
\newenvironment{case}[1][]{%
	\refstepcounter{case}%
	\refstepcounter{globalcase}% This creates unique anchors
	\par\smallskip%
	\noindent\textbf{Case \thecase\if\relax\detokenize{#1}\relax\else\ (#1)\fi.}\ %
}{%
	\par\smallskip%
}
\renewcommand{\theglobalcase}{\arabic{case}}

\AtBeginEnvironment{proof}{%
	\setcounter{case}{0}%
	\setcounter{step}{0}%
}

% 设置cleveref引用格式
\crefname{definition}{Definition}{Definitions}
\crefname{proposition}{Proposition}{Propositions}
\crefname{theorem}{Theorem}{Theorems}
\crefname{lemma}{Lemma}{Lemmas}
\crefname{corollary}{Corollary}{Corollaries}
\crefname{conjecture}{Conjecture}{Conjectures}
\crefname{question}{Question}{Questions}
\crefname{remark}{Remark}{Remarks}
\crefname{example}{Example}{Examples}
\crefname{proof}{Proof}{Proofs}
\crefname{claim}{Claim}{Claims}
\crefname{globalstep}{Step}{Steps} % Global unique step counter
\crefname{globalcase}{Case}{Cases} % Global unique case counter
\crefname{section}{Section}{Sections}
\crefname{subsection}{Subsection}{Subsections}
\crefname{mainthm}{Theorem}{Theorems}
